\fichatitulo{34 · ESCALA}{NeurIPS · 2022}{Training Compute-Optimal Large Language Models}
{\small Hoffmann et al. \textit{Training Compute-Optimal Large Language Models}. NeurIPS · 2022. \href{https://proceedings.neurips.cc/paper_files/paper/2022/file/c1e2faff6f588870935f114ebe04a3e5-Paper-Conference.pdf}{Fonte consultada}.\par}

\campo{Problema e objetivo}
Como distribuir um orçamento fixo entre parâmetros e tokens de treinamento?

\campo{Principal contribuição · síntese interpretativa}
Reavalia a alocação ótima de computação, enfatizando crescimento conjunto de modelo e dados.

\campo{Método / natureza da evidência}
Estuda mais de 400 modelos e apresenta Chinchilla, com 70 bilhões de parâmetros, treinado com mais dados que Gopher sob orçamento semelhante.

\campo{Resultado ou argumento principal}
Os resultados sustentam dimensionar conjuntamente parâmetros e tokens no regime estudado.

\campo{Excertos-chave · seleção literal}
\begin{itemize}
\excerto{1}{the model size and the number of training tokens should be scaled equally}{Resumo.}
\excerto{2}{under a given compute budget}{Resumo.}
\end{itemize}

\campo{Limitações e análise crítica}
Análise da ficha: a comparação depende de dados, treinamento e orçamento. Não implica que um gerador maior ou menor será superior no corpus da dissertação.

\campo{Aproveitamento na dissertação · proposta}
Fundamentação: comparar escala e eficiência; explicitar por que avaliação local continua necessária.

\registro{Fonte consultada nesta preparação: Resumo e introdução. Chave: \texttt{hoffmannEtAl2022computeOptimal}.}
